% !TeX program = xelatex
\documentclass[UTF8,11pt,a4paper,openany]{ctexrep}

\usepackage[a4paper,margin=2.15cm]{geometry}
\usepackage{amsmath,amssymb}
\usepackage{booktabs}
\usepackage{longtable}
\usepackage{array}
\usepackage{enumitem}
\usepackage{xcolor}
\usepackage{hyperref}
\usepackage{graphicx}
\usepackage{listings}
\usepackage[most]{tcolorbox}
\usepackage{tikz}
\usetikzlibrary{arrows.meta,positioning,fit,calc}

\definecolor{RosBlue}{HTML}{2563EB}
\definecolor{GoGreen}{HTML}{16A34A}
\definecolor{RssiOrange}{HTML}{EA580C}
\definecolor{WarnRed}{HTML}{B91C1C}
\definecolor{CodeBg}{HTML}{F6F8FA}
\definecolor{FrameGray}{HTML}{D0D7DE}
\definecolor{TextGray}{HTML}{57606A}

\hypersetup{
  colorlinks=true,
  linkcolor=RosBlue,
  citecolor=RosBlue,
  urlcolor=RosBlue,
  pdftitle={ROS 2 使用教程：面向 Unitree Go2 RSSI 机器狗项目},
  pdfauthor={Go2 RSSI Project}
}

\setlist[itemize]{leftmargin=2em,itemsep=0.2em,topsep=0.4em}
\setlist[enumerate]{leftmargin=2em,itemsep=0.2em,topsep=0.4em}
\renewcommand{\arraystretch}{1.18}
\setlength{\tabcolsep}{4pt}
\emergencystretch=2em
\newcolumntype{L}[1]{>{\raggedright\arraybackslash}p{#1}}
\urlstyle{same}

\lstdefinestyle{codestyle}{
  backgroundcolor=\color{CodeBg},
  basicstyle=\ttfamily\small,
  breaklines=true,
  breakatwhitespace=false,
  columns=fullflexible,
  keepspaces=true,
  frame=single,
  rulecolor=\color{FrameGray},
  xleftmargin=0.4em,
  xrightmargin=0.4em,
  aboveskip=0.8em,
  belowskip=0.8em,
  showstringspaces=false,
  tabsize=2
}
\lstdefinelanguage{yaml}{
  keywords={true,false,null,y,n},
  keywordstyle=\color{RosBlue}\bfseries,
  basicstyle=\ttfamily\small,
  sensitive=false,
  comment=[l]{\#},
  morecomment=[l]{---},
  morestring=[b]',
  morestring=[b]"
}
\lstdefinelanguage{bash}{
  keywords={cd,source,export,if,then,else,elif,fi,for,do,done,case,esac,function,exec,true,false,while,return},
  keywordstyle=\color{RosBlue}\bfseries,
  basicstyle=\ttfamily\small,
  sensitive=true,
  comment=[l]{\#},
  morestring=[b]',
  morestring=[b]"
}
\lstdefinelanguage{cpp}{
  keywords={class,public,private,void,int,double,bool,const,auto,return,if,else,while,for,namespace,using,new,static_cast,struct},
  keywordstyle=\color{RosBlue}\bfseries,
  basicstyle=\ttfamily\small,
  sensitive=true,
  comment=[l]{//},
  morecomment=[s]{/*}{*/},
  morestring=[b]"
}
\lstdefinelanguage{cmake}{
  keywords={cmake_minimum_required,project,if,endif,set,find_package,find_library,add_executable,target_include_directories,ament_target_dependencies,target_link_libraries,install,ament_package},
  keywordstyle=\color{RosBlue}\bfseries,
  basicstyle=\ttfamily\small,
  sensitive=false,
  comment=[l]{\#},
  morestring=[b]"
}
\lstset{style=codestyle}

\newtcolorbox{notebox}[1]{
  enhanced,
  breakable,
  colback=blue!3,
  colframe=RosBlue!65,
  title=#1,
  fonttitle=\bfseries,
  boxrule=0.6pt,
  arc=2mm,
  left=1mm,
  right=1mm,
  top=1mm,
  bottom=1mm
}

\newtcolorbox{warnbox}[1]{
  enhanced,
  breakable,
  colback=red!4,
  colframe=WarnRed!70,
  title=#1,
  fonttitle=\bfseries,
  boxrule=0.6pt,
  arc=2mm,
  left=1mm,
  right=1mm,
  top=1mm,
  bottom=1mm
}

\newtcolorbox{termdef}[1]{
  enhanced,
  breakable,
  colback=green!3,
  colframe=GoGreen!65,
  title=#1,
  fonttitle=\bfseries,
  boxrule=0.6pt,
  arc=2mm,
  left=1mm,
  right=1mm,
  top=1mm,
  bottom=1mm
}

\newcommand{\code}[1]{\nolinkurl{#1}}
\newcommand{\file}[1]{\nolinkurl{#1}}
\newcommand{\pkg}[1]{\nolinkurl{#1}}
\newcommand{\topic}[1]{\nolinkurl{#1}}
\newcommand{\rosnode}[1]{\nolinkurl{#1}}
\newcommand{\msg}[1]{\nolinkurl{#1}}
\newcommand{\param}[1]{\nolinkurl{#1}}
\newcommand{\actionname}[1]{\nolinkurl{#1}}

\title{
  \Huge ROS 2 使用教程\\
  \Large 面向 Unitree Go2 RSSI 机器狗项目的对照学习手册
}
\author{Go2 RSSI / ROS 2 Integration Notes}
\date{2026-06-15}

\begin{document}

\maketitle
\tableofcontents

\chapter*{阅读方式}
\addcontentsline{toc}{chapter}{阅读方式}

这份教程不是泛泛介绍 ROS 2，而是围绕本项目已经存在的代码来讲。项目当前的 ROS 2 核心包是 \pkg{go2_nav2_bridge}，它把 Unitree Go2 的 DDS 数据转换成标准 ROS 2 话题，并把 ROS 2 的速度命令转换成 \code{SportClient::Move}。因此本文的主线是：

\begin{center}
  \textbf{先理解 ROS 2 的基本通信模型，再把每个概念对应到本项目的 Go2 bridge、RSSI/IMU 采集、Nav2 后端和直线控制器。}
\end{center}

建议按下面顺序读：

\begin{enumerate}
  \item 第 \ref{chap:big-picture} 章先建立 ROS 2 在本项目里的总体图。
  \item 第 \ref{chap:core} 章学习 node、topic、message、parameter、launch、TF、QoS、action 等基础概念。
  \item 第 \ref{chap:project-package} 章对照 \pkg{go2_nav2_bridge} 的文件结构和 CMake/package 配置。
  \item 第 \ref{chap:bridge} 章细看 Go2 bridge 怎样发布 \topic{/odom}、\topic{/go2/imu}、\topic{/go2/path_s}、\topic{/go2/pointcloud}，以及怎样接收 \topic{/cmd_vel}。
  \item 第 \ref{chap:run-debug} 章学习如何构建、运行、看话题、查 TF、查参数和定位问题。
\end{enumerate}

\begin{notebox}{本文约定}
本文中的“本项目”指当前仓库中的 Go2 RSSI 路径进度项目。ROS 2 主要运行在外部电脑上，Go2 侧通过 Unitree SDK2/DDS 提供真实运动状态、IMU、点云和运动控制接口。RSSI/IMU 融合算法的主输出是路径进度 \(s\)，ROS 2 的任务是把机器人运动、传感器、采集脚本和可选 Nav2 后端组织起来。
\end{notebox}

\chapter{总体视角：ROS 2 在本项目里解决什么问题}
\label{chap:big-picture}

\section{一句话理解 ROS 2}

ROS 2 可以先简单理解成“机器人软件的通信和组织框架”。它不直接等于控制算法，也不直接等于导航算法，而是负责让很多独立程序可靠地互相通信：

\begin{itemize}
  \item Go2 bridge 负责和 Unitree SDK2/DDS 打交道。
  \item 直线控制器或 Nav2 负责产生速度命令。
  \item RSSI/IMU 采集程序负责保存实验数据。
  \item RViz、命令行工具、诊断脚本负责观察系统状态。
\end{itemize}

这些程序不需要写在同一个进程里。它们通过 ROS 2 的 topic、service、action、parameter、TF 等机制连接。

\section{本项目的数据和控制闭环}

\begin{figure}[htbp]
\centering
\resizebox{0.98\textwidth}{!}{
\begin{tikzpicture}[
  font=\small,
  box/.style={draw, rounded corners=2mm, align=center, minimum width=31mm, minimum height=12mm, inner sep=2mm},
  ros/.style={box, fill=blue!5, draw=RosBlue!70},
  go/.style={box, fill=green!5, draw=GoGreen!70},
  rssi/.style={box, fill=orange!6, draw=RssiOrange!75},
  arrow/.style={-Latex, thick}
]
  \node[go] (go2) {Unitree Go2\\DDS / SportClient};
  \node[ros, right=22mm of go2] (bridge) {\rosnode{go2_nav2_bridge}\\ROS 2 硬件桥};
  \node[ros, above right=14mm and 22mm of bridge] (nav2) {Nav2 或\\直线控制器};
  \node[ros, below right=14mm and 22mm of bridge] (topics) {ROS 2 话题\\odom / imu / path\_s / pointcloud};
  \node[rssi, right=26mm of topics] (fusion) {RSSI + IMU/odom\\融合算法};
  \node[rssi, below=14mm of topics] (collector) {采集脚本\\run directory CSV/log};

  \draw[arrow] (go2) -- node[above, align=center]{sport state\\point cloud} (bridge);
  \draw[arrow] (bridge) -- node[above, align=center]{\topic{/odom}\\\topic{/go2/imu}\\\topic{/go2/path_s}} (topics);
  \draw[arrow] (topics) -- node[above]{同步窗口} (fusion);
  \draw[arrow] (nav2) -- node[above, align=center]{\topic{/cmd_vel}\\或 \topic{/cmd_vel_nav}} (bridge);
  \draw[arrow] (bridge) -- node[below]{\code{SportClient::Move}} (go2);
  \draw[arrow] (topics) -- (collector);
  \draw[arrow] (collector) -- node[right]{RSSI/IMU/odom 记录} (fusion);
\end{tikzpicture}}
\caption{本项目中的 ROS 2 角色：桥接真实 Go2、组织运动命令、发布标准传感器话题，并为 RSSI/IMU 融合算法提供时间对齐的数据。}
\label{fig:project-ros2-flow}
\end{figure}

从图 \ref{fig:project-ros2-flow} 可以看到，本项目里 ROS 2 最核心的价值有三个：

\begin{enumerate}
  \item \textbf{统一接口。} Go2 原始 DDS 数据被转换成 \msg{nav_msgs/msg/Odometry}、\msg{sensor_msgs/msg/Imu}、\msg{sensor_msgs/msg/PointCloud2} 等标准消息。
  \item \textbf{解耦模块。} 直线控制器、Nav2、诊断脚本、采集脚本都通过话题或 action 接入，不需要直接依赖 Unitree SDK2。
  \item \textbf{可观测和可调试。} 可以用 \code{ros2 topic echo}、\code{ros2 topic hz}、\code{ros2 param get}、RViz 等工具检查系统。
\end{enumerate}

\section{本项目当前 ROS 2 文件总表}

\begin{longtable}{L{0.30\textwidth} L{0.25\textwidth} L{0.36\textwidth}}
\toprule
\textbf{文件} & \textbf{ROS 2 角色} & \textbf{在本项目中的作用} \\
\midrule
\endhead
\file{ros2_ws/src/go2_nav2_bridge/package.xml}
& package 元数据
& 声明包名、依赖、运行时组件。ROS 2 用它识别 \pkg{go2_nav2_bridge}。\\
\file{ros2_ws/src/go2_nav2_bridge/CMakeLists.txt}
& 构建规则
& 用 \code{ament_cmake} 编译三个 C++ 可执行节点，并链接 Unitree SDK2。\\
\file{src/go2_nav2_bridge.cpp}
& ROS 2 node
& 订阅 Unitree DDS，发布 \topic{/odom}、\topic{/go2/imu}、\topic{/go2/path_s}、\topic{/go2/pointcloud}，订阅 \topic{/cmd_vel} 后调用 \code{SportClient::Move}。\\
\file{src/go2_direct_straight_controller.cpp}
& ROS 2 node
& 直接发布 \topic{/cmd_vel}，按 \topic{/go2/path_s}、\topic{/odom} 或时间判断何时停止。\\
\file{src/nav2_straight_path_client.cpp}
& action client
& 向 Nav2 发送 \actionname{FollowPath}、\actionname{ComputePathToPose} 或 \actionname{NavigateToPose} 目标。\\
\file{launch/go2_nav2_straight.launch.py}
& launch 文件
& 一次性启动 Go2 bridge、Nav2 controller、planner、BT navigator、velocity smoother、collision monitor 等节点。\\
\file{params/go2_nav2_minimal_controller_params.yaml}
& 参数文件
& direct/minimal 场景参数，适合无障碍直线 RSSI 采集。\\
\file{params/go2_nav2_straight_params.yaml}
& 参数文件
& obstacle/Nav2 场景参数，打开点云、costmap、velocity smoother 和 collision monitor。\\
\file{scripts/go2_collect_straight_line_nav2.sh}
& 实验编排脚本
& 启动 RSSI、IMU、Go2 bridge、direct 或 Nav2 后端，并写入 run directory。\\
\file{scripts/go2_diagnose_nav2_perception.sh}
& 诊断脚本
& 不让机器人运动，只启动感知链路，检查点云、costmap 和 collision monitor 前置条件。\\
\bottomrule
\end{longtable}

\chapter{ROS 2 核心概念：用本项目来理解}
\label{chap:core}

\section{Workspace：工作空间}

\begin{termdef}{Workspace}
ROS 2 workspace 是一组 package 的开发目录。本项目的 ROS 2 workspace 是 \file{code/ros2_ws}。
\end{termdef}

本项目的结构是：

\begin{lstlisting}[language=bash]
code/
  ros2_ws/
    src/
      go2_nav2_bridge/
        package.xml
        CMakeLists.txt
        launch/
        params/
        src/
\end{lstlisting}

常用构建命令：

\begin{lstlisting}[language=bash]
cd /home/luping/桌面/RSSI/RSSI/code/ros2_ws
source /opt/ros/humble/setup.bash
colcon build --packages-select go2_nav2_bridge
source install/setup.bash
\end{lstlisting}

\begin{warnbox}{为什么一定要 source}
\code{source /opt/ros/humble/setup.bash} 让 shell 知道 ROS 2 基础环境在哪里；\code{source install/setup.bash} 让 shell 知道本项目编译出的 \pkg{go2_nav2_bridge} 在哪里。没有 source，常见现象是 \code{ros2 pkg prefix go2_nav2_bridge} 找不到包，或者 \code{ros2 run go2_nav2_bridge ...} 找不到可执行文件。
\end{warnbox}

\section{Package：功能包}

\begin{termdef}{Package}
package 是 ROS 2 的最小组织单元。一个 package 可以包含节点源码、launch 文件、参数文件、消息定义、测试和安装规则。
\end{termdef}

本项目当前的 ROS 2 package 是 \pkg{go2_nav2_bridge}。它的 \file{package.xml} 声明了依赖：

\begin{lstlisting}[language=xml]
<depend>rclcpp</depend>
<depend>geometry_msgs</depend>
<depend>nav_msgs</depend>
<depend>sensor_msgs</depend>
<depend>tf2_ros</depend>
<depend>nav2_msgs</depend>
\end{lstlisting}

这些依赖的含义：

\begin{longtable}{L{0.22\textwidth} L{0.68\textwidth}}
\toprule
\textbf{依赖} & \textbf{作用} \\
\midrule
\endhead
\pkg{rclcpp} & C++ 节点 API，用来创建 node、publisher、subscriber、timer、parameter。\\
\pkg{rclcpp_action} & C++ action API，用来和 Nav2 的 \actionname{FollowPath}、\actionname{NavigateToPose} 交互。\\
\pkg{geometry_msgs} & 提供 \msg{Twist}、\msg{PoseStamped}、\msg{TransformStamped} 等几何消息。\\
\pkg{nav_msgs} & 提供 \msg{Odometry}、\msg{Path} 等导航消息。\\
\pkg{sensor_msgs} & 提供 \msg{Imu}、\msg{PointCloud2} 等传感器消息。\\
\pkg{std_msgs} & 提供 \msg{Float32} 等简单消息，本项目用它发布 \topic{/go2/path_s}。\\
\pkg{tf2_ros} & 发布和查询坐标变换，例如 \code{odom -> base_link}。\\
\pkg{nav2_msgs} & Nav2 action 和消息类型。\\
\bottomrule
\end{longtable}

\section{Node：节点}

\begin{termdef}{Node}
node 是一个正在运行的 ROS 2 程序。一个节点通常负责一件清晰的事情，例如桥接硬件、发布速度、发送导航目标、诊断点云。
\end{termdef}

本项目最重要的节点如下：

\begin{longtable}{L{0.30\textwidth} L{0.24\textwidth} L{0.36\textwidth}}
\toprule
\textbf{节点名} & \textbf{来源} & \textbf{职责} \\
\midrule
\endhead
\rosnode{go2_nav2_bridge} & 本项目 C++ & 硬件桥。把 Go2 状态转换成 ROS 2 话题，把 \topic{/cmd_vel} 转成 Go2 运动命令。\\
\rosnode{go2_direct_straight_controller} & 本项目 C++ & 直线采集控制器。发固定前进速度，按距离或时间停止。\\
\rosnode{nav2_straight_path_client} & 本项目 C++ & Nav2 客户端。发送直线路径或目标点给 Nav2。\\
\rosnode{controller_server} & Nav2 & 根据路径和 costmap 产生局部速度命令。\\
\rosnode{velocity_smoother} & Nav2 & 平滑速度命令，限制加速度和减速度。\\
\rosnode{collision_monitor} & Nav2 & 根据点云安全区修改或阻断速度命令。\\
\rosnode{go2_nav2_perception_diag} & 本项目 Python & 订阅点云和 costmap，输出感知链路统计。\\
\bottomrule
\end{longtable}

查看当前节点：

\begin{lstlisting}[language=bash]
ros2 node list
ros2 node info /go2_nav2_bridge
\end{lstlisting}

\section{Topic 和 Message：话题与消息}

\begin{termdef}{Topic}
topic 是连续数据流。publisher 往 topic 发消息，subscriber 从 topic 收消息。topic 适合传感器、速度、状态这类“不断更新”的数据。
\end{termdef}

\begin{termdef}{Message}
message 是 topic 上每一帧数据的结构。例如 \topic{/cmd_vel} 的消息类型是 \msg{geometry_msgs/msg/Twist}，里面有线速度和角速度。
\end{termdef}

本项目的核心 topic：

\begin{longtable}{L{0.22\textwidth} L{0.28\textwidth} L{0.40\textwidth}}
\toprule
\textbf{Topic} & \textbf{消息类型} & \textbf{含义} \\
\midrule
\endhead
\topic{/cmd_vel} & \msg{geometry_msgs/msg/Twist} & 最终发给 Go2 bridge 的速度命令。bridge 会调用 \code{SportClient::Move(vx, vy, wz)}。\\
\topic{/cmd_vel_nav} & \msg{geometry_msgs/msg/Twist} & Nav2 controller 输出的原始速度，之后进入 velocity smoother。\\
\topic{/cmd_vel_smoothed} & \msg{geometry_msgs/msg/Twist} & velocity smoother 输出的平滑速度，之后进入 collision monitor。\\
\topic{/odom} & \msg{nav_msgs/msg/Odometry} & Go2 当前位姿和速度。来源是 Unitree \code{SportModeState}。\\
\topic{/tf} & \msg{tf2_msgs/msg/TFMessage} & 坐标变换流，本项目主要发布 \code{odom -> base_link}。\\
\topic{/go2/imu} & \msg{sensor_msgs/msg/Imu} & Go2 IMU 姿态、角速度、线加速度。来源是 \code{SportModeState.imu_state()}。\\
\topic{/go2/path_s} & \msg{std_msgs/msg/Float32} & bridge 根据 Go2 position 增量累计出的路径进度。\\
\topic{/go2/pointcloud} & \msg{sensor_msgs/msg/PointCloud2} & Go2 点云，供 Nav2 costmap 和 collision monitor 使用。\\
\topic{/local_costmap/costmap} & \msg{nav_msgs/msg/OccupancyGrid} & Nav2 局部代价地图，诊断脚本会读取它。\\
\bottomrule
\end{longtable}

常用命令：

\begin{lstlisting}[language=bash]
ros2 topic list
ros2 topic info /odom
ros2 topic echo /go2/path_s
ros2 topic hz /go2/imu
ros2 interface show geometry_msgs/msg/Twist
\end{lstlisting}

\section{Publisher 和 Subscriber：发布与订阅}

在 \file{go2_nav2_bridge.cpp} 里，bridge 创建 publisher：

\begin{lstlisting}[language=cpp]
odom_pub_ = create_publisher<nav_msgs::msg::Odometry>(odom_topic_, rclcpp::SensorDataQoS());
imu_pub_ = create_publisher<sensor_msgs::msg::Imu>(imu_topic_, rclcpp::SensorDataQoS());
path_s_pub_ = create_publisher<std_msgs::msg::Float32>(path_s_topic_, 10);
point_cloud_pub_ = create_publisher<sensor_msgs::msg::PointCloud2>(
  ros_point_cloud_topic_, rclcpp::SensorDataQoS());
\end{lstlisting}

也创建 subscriber：

\begin{lstlisting}[language=cpp]
cmd_vel_sub_ = create_subscription<geometry_msgs::msg::Twist>(
  cmd_vel_topic_, 10,
  std::bind(&Go2Nav2Bridge::onTwist, this, std::placeholders::_1));
\end{lstlisting}

直白理解：

\begin{itemize}
  \item publisher 是“往外广播数据”。
  \item subscriber 是“监听别人发来的数据”。
  \item bridge 既是 publisher，也是 subscriber，因为它既发布 Go2 状态，也接收运动命令。
\end{itemize}

\section{Parameter：参数}

\begin{termdef}{Parameter}
parameter 是节点启动时或运行时可配置的变量。它让同一个程序在不同实验场景下使用不同配置，而不需要改 C++ 源码。
\end{termdef}

\file{go2_nav2_bridge.cpp} 中有大量 \code{declare_parameter}：

\begin{lstlisting}[language=cpp]
network_interface_ = declare_parameter<std::string>("network_interface", "");
cmd_vel_topic_ = declare_parameter<std::string>("cmd_vel_topic", "/cmd_vel");
odom_topic_ = declare_parameter<std::string>("odom_topic", "/odom");
imu_topic_ = declare_parameter<std::string>("imu_topic", "/go2/imu");
publish_point_cloud_ = declare_parameter<bool>("publish_point_cloud", true);
cmd_timeout_sec_ = declare_parameter<double>("cmd_timeout_sec", 0.5);
max_vx_mps_ = declare_parameter<double>("max_vx_mps", 0.35);
\end{lstlisting}

这些参数在 YAML 中配置：

\begin{lstlisting}[language=yaml]
go2_nav2_bridge:
  ros__parameters:
    network_interface: eth0
    cmd_vel_topic: /cmd_vel
    odom_topic: /odom
    imu_topic: /go2/imu
    publish_point_cloud: true
    cmd_timeout_sec: 0.6
    max_vx_mps: 0.12
\end{lstlisting}

查询参数：

\begin{lstlisting}[language=bash]
ros2 param list /go2_nav2_bridge
ros2 param get /go2_nav2_bridge max_vx_mps
\end{lstlisting}

\section{Launch：一键启动多个节点}

\begin{termdef}{Launch}
launch 文件是 ROS 2 的启动编排。它能一次启动多个节点，并给它们传参数、重映射话题、设置条件。
\end{termdef}

本项目的 \file{go2_nav2_straight.launch.py} 做了这些事：

\begin{itemize}
  \item 永远启动 \rosnode{go2_nav2_bridge}。
  \item 当 \param{move_backend=nav2} 时启动 Nav2 controller、velocity smoother、collision monitor。
  \item 当 \param{use_planner=true} 或 \param{use_bt_navigator=true} 时启动 planner。
  \item 当 \param{use_bt_navigator=true} 时启动 behavior server 和 BT navigator。
  \item 启动 lifecycle manager，使 Nav2 管理节点进入 active 状态。
\end{itemize}

启动示例：

\begin{lstlisting}[language=bash]
ros2 launch go2_nav2_bridge go2_nav2_straight.launch.py \
  params_file:=/path/to/go2_nav2_straight_params.yaml \
  move_backend:=nav2 \
  use_bt_navigator:=false \
  use_planner:=false \
  network_interface:=enp5s0
\end{lstlisting}

\section{Remapping：话题重映射}

Nav2 controller 默认输出 \topic{cmd_vel}。本项目在 launch 里把它重映射成 \topic{/cmd_vel_nav}：

\begin{lstlisting}[language=python]
Node(
  package="nav2_controller",
  executable="controller_server",
  name="controller_server",
  remappings=[("cmd_vel", "/cmd_vel_nav")],
)
\end{lstlisting}

这样做的原因是要把速度链路拆开：

\begin{center}
\topic{/cmd_vel_nav} \(\rightarrow\) \rosnode{velocity_smoother} \(\rightarrow\)
\topic{/cmd_vel_smoothed} \(\rightarrow\) \rosnode{collision_monitor} \(\rightarrow\)
\topic{/cmd_vel} \(\rightarrow\) \rosnode{go2_nav2_bridge}
\end{center}

也就是说，\topic{/cmd_vel} 是真正进入 Go2 bridge 的最终速度命令。

\section{TF：坐标变换}

\begin{termdef}{TF}
TF 是 ROS 2 中维护坐标系关系的系统。机器人导航至少需要知道 \code{odom} 坐标系和 \code{base_link} 坐标系之间的关系。
\end{termdef}

本项目 bridge 发布：

\begin{center}
  \code{odom -> base_link}
\end{center}

\file{go2_nav2_bridge.cpp} 中相关逻辑是：

\begin{lstlisting}[language=cpp]
tf_msg.header.frame_id = odom_frame_;
tf_msg.child_frame_id = base_frame_;
tf_msg.transform.translation.x = odom.pose.pose.position.x;
tf_msg.transform.translation.y = odom.pose.pose.position.y;
tf_msg.transform.rotation = odom.pose.pose.orientation;
tf_broadcaster_->sendTransform(tf_msg);
\end{lstlisting}

检查 TF：

\begin{lstlisting}[language=bash]
ros2 run tf2_ros tf2_echo odom base_link
\end{lstlisting}

\begin{warnbox}{TF 出问题时会怎样}
Nav2 controller、planner、costmap 都依赖 TF。如果 \code{odom -> base_link} 不存在、时间戳太旧、frame 名不一致，常见现象是 Nav2 active 了但不出速度，或者日志里出现 transform timeout。
\end{warnbox}

\section{QoS：通信质量配置}

\begin{termdef}{QoS}
QoS 决定 ROS 2 通信的可靠性、缓存深度、是否保留上一帧等策略。传感器数据通常用低延迟策略，地图类数据常用可靠和 transient local 策略。
\end{termdef}

本项目中：

\begin{itemize}
  \item \topic{/odom}、\topic{/go2/imu}、\topic{/go2/pointcloud} 使用 \code{rclcpp::SensorDataQoS()}，适合高频传感器。
  \item \topic{/go2/path_s} 使用普通 depth 10，因为它是轻量路径进度值。
  \item 诊断脚本订阅 \topic{/go2/pointcloud} 时使用 best effort，订阅 \topic{/local_costmap/costmap} 时使用 reliable + transient local。
\end{itemize}

如果 \code{ros2 topic echo} 看不到某些传感器数据，但节点确实在发布，可能是 QoS 不匹配。可用：

\begin{lstlisting}[language=bash]
ros2 topic info /go2/pointcloud --verbose
\end{lstlisting}

\section{Action：长任务接口}

\begin{termdef}{Action}
action 用于“会持续一段时间、需要反馈、可以取消”的任务。例如导航到目标点、跟随路径。
\end{termdef}

本项目的 \rosnode{nav2_straight_path_client} 使用三个 Nav2 action：

\begin{longtable}{L{0.28\textwidth} L{0.56\textwidth}}
\toprule
\textbf{Action} & \textbf{本项目用途} \\
\midrule
\endhead
\actionname{FollowPath} & 给 Nav2 controller 一条路径，让 controller 产生局部速度。\\
\actionname{ComputePathToPose} & 请求 planner 从当前位置规划到目标点。\\
\actionname{NavigateToPose} & 交给 BT navigator 完成目标导航，可包含重规划和行为树逻辑。\\
\bottomrule
\end{longtable}

初学时可以这样区分：

\begin{itemize}
  \item topic 像“直播数据流”：IMU、odom、cmd\_vel。
  \item service 像“一问一答”：请求一次，返回一次。
  \item action 像“执行任务”：发送目标，期间有反馈，最后有结果，还可以取消。
\end{itemize}

\section{Lifecycle：Nav2 管理节点状态}

Nav2 很多节点不是普通 node，而是 lifecycle node。它们有状态：unconfigured、inactive、active、finalized。只有 active 后才真正工作。

本项目 launch 使用 \rosnode{lifecycle_manager} 自动激活 Nav2 节点。采集脚本里也会等待日志出现：

\begin{lstlisting}[language=bash]
Managed nodes are active
\end{lstlisting}

命令行检查：

\begin{lstlisting}[language=bash]
ros2 lifecycle get /controller_server
\end{lstlisting}

\chapter{本项目 ROS 2 Package 逐文件对照}
\label{chap:project-package}

\section{package.xml：告诉 ROS 2 这个包是什么}

\file{package.xml} 里最关键的是三类信息：

\begin{itemize}
  \item 包名：\pkg{go2_nav2_bridge}。
  \item 构建工具：\pkg{ament_cmake}。
  \item 编译和运行依赖：\pkg{rclcpp}、\pkg{nav_msgs}、\pkg{sensor_msgs}、\pkg{tf2_ros}、\pkg{nav2_msgs}、Nav2 相关包。
\end{itemize}

本项目还声明了 Nav2 运行依赖：

\begin{lstlisting}[language=xml]
<exec_depend>nav2_controller</exec_depend>
<exec_depend>nav2_costmap_2d</exec_depend>
<exec_depend>nav2_velocity_smoother</exec_depend>
<exec_depend>nav2_collision_monitor</exec_depend>
<exec_depend>nav2_lifecycle_manager</exec_depend>
\end{lstlisting}

这表示：编译本项目 C++ 代码不一定需要直接链接这些 Nav2 进程，但运行 launch 时需要它们存在。

\section{CMakeLists.txt：告诉 colcon 怎么编译}

\file{CMakeLists.txt} 中定义了三个可执行节点：

\begin{lstlisting}[language=cmake]
add_executable(go2_nav2_bridge src/go2_nav2_bridge.cpp)
add_executable(nav2_straight_path_client src/nav2_straight_path_client.cpp)
add_executable(go2_direct_straight_controller src/go2_direct_straight_controller.cpp)
\end{lstlisting}

然后安装：

\begin{lstlisting}[language=cmake]
install(TARGETS
  go2_nav2_bridge
  nav2_straight_path_client
  go2_direct_straight_controller
  DESTINATION lib/${PROJECT_NAME})

install(DIRECTORY launch params
  DESTINATION share/${PROJECT_NAME})
\end{lstlisting}

这两段决定了：

\begin{itemize}
  \item \code{ros2 run go2_nav2_bridge go2_nav2_bridge} 可以找到可执行文件。
  \item \code{ros2 launch go2_nav2_bridge go2_nav2_straight.launch.py} 可以找到 launch 文件。
  \item \code{ros2 pkg prefix go2_nav2_bridge} 下可以找到安装后的 params 文件。
\end{itemize}

\section{Unitree SDK2 链接}

本项目 bridge 需要链接 Unitree SDK2：

\begin{lstlisting}[language=cmake]
find_library(UNITREE_SDK2_LIB
  NAMES unitree_sdk2
  PATHS "${UNITREE_SDK2_DIR}/lib/${UNITREE_SDK2_ARCH}"
  NO_DEFAULT_PATH)
\end{lstlisting}

如果 SDK2 不在默认位置，可以这样构建：

\begin{lstlisting}[language=bash]
colcon build --packages-select go2_nav2_bridge \
  --cmake-args -DUNITREE_SDK2_DIR=/path/to/unitree_sdk2
\end{lstlisting}

运行时还要保证动态库路径正确。采集脚本会把：

\begin{lstlisting}[language=bash]
code/unitree_sdk2/thirdparty/lib/$(uname -m)
\end{lstlisting}

加入 \code{LD_LIBRARY_PATH}。

\chapter{Go2 Bridge：本项目最关键的 ROS 2 节点}
\label{chap:bridge}

\section{Bridge 的输入输出}

\rosnode{go2_nav2_bridge} 是本项目 ROS 2 链路的中心。它一边连接 Unitree DDS/SportClient，一边连接 ROS 2。

\begin{longtable}{L{0.23\textwidth} L{0.18\textwidth} L{0.48\textwidth}}
\toprule
\textbf{方向} & \textbf{接口} & \textbf{说明} \\
\midrule
\endhead
Go2 \(\rightarrow\) ROS 2 & \code{rt/sportmodestate} & 读取 Go2 position、velocity、rpy、gyro、acc，转换成 odom、imu、path\_s、TF。\\
Go2 \(\rightarrow\) ROS 2 & \code{rt/utlidar/cloud} & 读取 Unitree 点云，转换成 \msg{sensor_msgs/msg/PointCloud2}，可做外参变换和自过滤。\\
ROS 2 \(\rightarrow\) Go2 & \topic{/cmd_vel} & 接收线速度和角速度，限幅后调用 \code{SportClient::Move}。\\
\bottomrule
\end{longtable}

\section{速度命令如何进入 Go2}

当 bridge 收到 \topic{/cmd_vel}：

\begin{lstlisting}[language=cpp]
void onTwist(const geometry_msgs::msg::Twist::SharedPtr msg)
{
  sendVelocity(msg->linear.x, msg->linear.y, msg->angular.z);
}
\end{lstlisting}

\code{sendVelocity} 会做三件重要的事：

\begin{enumerate}
  \item 把非有限值变成 0，避免 NaN/inf 进入机器人。
  \item 用 \param{max_vx_mps}、\param{max_vy_mps}、\param{max_wz_radps} 对速度限幅。
  \item 调用 \code{sport_client_->Move(vx, vy, wz)}。
\end{enumerate}

简化逻辑：

\begin{lstlisting}[language=cpp]
safe_vx = clamp(vx, -max_vx_mps_, max_vx_mps_);
safe_vy = clamp(vy, -max_vy_mps_, max_vy_mps_);
safe_wz = clamp(wz, -max_wz_radps_, max_wz_radps_);
sport_client_->Move(safe_vx, safe_vy, safe_wz);
\end{lstlisting}

\begin{warnbox}{速度限幅是安全边界}
Nav2、直线控制器、命令行都可能发 \topic{/cmd_vel}。真正进入 Go2 前，bridge 的 \param{max_vx_mps}、\param{max_vy_mps}、\param{max_wz_radps} 是最后一道软件限幅。实验时不要只看 controller 参数，也要看 bridge 参数。
\end{warnbox}

\section{watchdog：为什么机器人会自动停}

bridge 中有 watchdog timer。若持续一段时间没有新的非零速度命令，就发送停止：

\begin{lstlisting}[language=cpp]
if (age > cmd_timeout_sec_) {
  sendStop();
}
\end{lstlisting}

这对应参数：

\begin{lstlisting}[language=yaml]
cmd_timeout_sec: 0.6
control_watchdog_hz: 20.0
\end{lstlisting}

直白理解：ROS 2 速度命令必须持续发布。如果上游节点崩了、网络断了、Nav2 卡了，bridge 不会让 Go2 一直按旧速度走。

\section{odom 是怎么来的}

\topic{/odom} 不是 RSSI 算出来的，也不是 Nav2 算出来的。本项目的 \topic{/odom} 来自 Go2 的 \code{SportModeState}：

\begin{lstlisting}[language=cpp]
const auto & position = state.position();
const auto & velocity = state.velocity();
const auto & rpy = state.imu_state().rpy();
\end{lstlisting}

bridge 把这些字段填入 \msg{nav_msgs/msg/Odometry}：

\begin{lstlisting}[language=cpp]
odom.pose.pose.position.x = position[0];
odom.pose.pose.position.y = position[1];
odom.pose.pose.position.z = position[2];
odom.twist.twist.linear.x = velocity[0];
odom.twist.twist.linear.y = velocity[1];
odom.twist.twist.angular.z = state.yaw_speed();
\end{lstlisting}

所以本项目里的 odom 可以理解为：

\begin{center}
  \textbf{Go2 底层状态估计给出的短期位姿和速度，bridge 把它包装成标准 ROS 2 odometry。}
\end{center}

它适合做短时连续性、运动约束、Nav2 controller 输入和实验参考；不应把它当成 RSSI 研究算法的主定位结果。

\section{/go2/path\_s 是怎么来的}

bridge 每次收到 Go2 sport state，会比较当前 position 和上一帧 position：

\begin{lstlisting}[language=cpp]
const double dx = position[0] - last_x_;
const double dy = position[1] - last_y_;
const double ds = std::hypot(dx, dy);
path_s_m_ += ds;
\end{lstlisting}

然后发布：

\begin{lstlisting}[language=cpp]
std_msgs::msg::Float32 path_s;
path_s.data = static_cast<float>(path_s_m_);
path_s_pub_->publish(path_s);
\end{lstlisting}

\topic{/go2/path_s} 的含义是“从本次 bridge 启动后累计走过的路径长度”。它不是全局地图上的真实 \(s\)，但在直线采集时很有用，因为它可以告诉直线控制器“这一趟已经走了多少米”。

\section{IMU 是怎么发布的}

Go2 sport state 中的 IMU 字段被转成标准 \msg{sensor_msgs/msg/Imu}：

\begin{lstlisting}[language=cpp]
imu.orientation = odom.pose.pose.orientation;
imu.angular_velocity.x = gyro[0];
imu.angular_velocity.y = gyro[1];
imu.angular_velocity.z = gyro[2];
imu.linear_acceleration.x = acc[0];
imu.linear_acceleration.y = acc[1];
imu.linear_acceleration.z = acc[2];
imu_pub_->publish(imu);
\end{lstlisting}

融合算法使用 IMU/odom 的方式可以简单理解为：RSSI 给当前位置观测，IMU/odom 给短时间内“应该移动了多少”的运动约束，最后融合出预测 \(s\)。

\section{点云处理：外参、过滤和新鲜度检查}

当 \param{publish_point_cloud=true} 时，bridge 会把 Unitree 点云发布成 \topic{/go2/pointcloud}。中间会做三件事：

\begin{enumerate}
  \item 根据 \param{point_cloud_xyz_offset} 和 \param{point_cloud_rpy_offset} 做坐标变换。
  \item 根据 \param{self_filter_boxes} 删除机器人自身附近的点。
  \item 更新点云时间戳，给 \param{require_fresh_point_cloud_for_motion} 使用。
\end{enumerate}

如果打开：

\begin{lstlisting}[language=yaml]
require_fresh_point_cloud_for_motion: true
point_cloud_motion_timeout_sec: 0.8
\end{lstlisting}

那么点云太久没更新时，bridge 会阻止非零速度命令。这适合 obstacle/Nav2 测试，但不适合没有点云的 minimal 直线采集。

\chapter{直线控制器与 Nav2 客户端}

\section{直线控制器：最小闭环}

\rosnode{go2_direct_straight_controller} 是本项目推荐的无障碍直线采集后端。它不依赖 Nav2 controller，只做简单稳定的事情：

\begin{enumerate}
  \item 订阅 \topic{/go2/path_s} 和 \topic{/odom}。
  \item 按 \param{speed_mps} 发布 \topic{/cmd_vel}。
  \item 根据 \param{distance_source} 判断是否到达目标距离。
  \item 到达后发布 0 速度并退出。
\end{enumerate}

核心参数：

\begin{longtable}{L{0.24\textwidth} L{0.18\textwidth} L{0.48\textwidth}}
\toprule
\textbf{参数} & \textbf{常用值} & \textbf{含义} \\
\midrule
\endhead
\param{distance_m} & \code{1.0} & 本次直线采集目标距离。\\
\param{speed_mps} & \code{0.20} & 前进速度。源码中会限制在 \code{0.03--0.35}。\\
\param{max_runtime_sec} & \code{30.0} & 超时保护，避免一直走。\\
\param{control_hz} & \code{20.0} & 发布速度命令频率。\\
\param{distance_source} & \code{path_s} & 停止依据，可选 \code{path_s}、\code{odom}、\code{time}。\\
\param{cmd_vel_topic} & \topic{/cmd_vel} & 输出速度话题。\\
\bottomrule
\end{longtable}

运行示例：

\begin{lstlisting}[language=bash]
ros2 run go2_nav2_bridge go2_direct_straight_controller \
  --ros-args \
  --params-file /path/to/go2_nav2_minimal_controller_params.yaml \
  -p distance_m:=1.0 \
  -p speed_mps:=0.20 \
  -p distance_source:=path_s
\end{lstlisting}

\section{Nav2 客户端：把目标交给 Nav2}

\rosnode{nav2_straight_path_client} 是本项目的 Nav2 action 客户端。它的工作流程是：

\begin{enumerate}
  \item 通过 TF 查询当前 \code{odom -> base_link}。
  \item 根据当前位置生成前方 \param{distance_m} 的目标或路径。
  \item 如果 \param{use_planner=true}，请求 planner 生成路径。
  \item 如果 \param{use_navigate_to_pose=true}，走 BT navigator。
  \item 否则直接发送 \actionname{FollowPath}。
\end{enumerate}

常用参数：

\begin{longtable}{L{0.27\textwidth} L{0.20\textwidth} L{0.43\textwidth}}
\toprule
\textbf{参数} & \textbf{常用值} & \textbf{含义} \\
\midrule
\endhead
\param{distance_m} & \code{0.5--1.0} & 目标距离。\\
\param{path_resolution_m} & \code{0.10} & 自生成直线路径的点间距。\\
\param{odom_frame} / \param{base_frame} & \code{odom} / \code{base_link} & 查询 TF 的坐标系。\\
\param{controller_id} & \code{FollowPath} & controller server 中的插件 ID。\\
\param{action_name} & \code{follow_path} & FollowPath action 名称。\\
\param{use_planner} & \code{false} 或 \code{true} & 是否先请求 \actionname{ComputePathToPose}。\\
\param{use_navigate_to_pose} & \code{false} 或 \code{true} & 是否使用 BT navigator。\\
\param{finish_on_path_s} & \code{false} & 是否按 \topic{/go2/path_s} 到达距离后取消 action。\\
\bottomrule
\end{longtable}

\chapter{Launch 与参数文件：本项目怎样一键启动}

\section{Launch 参数}

\file{go2_nav2_straight.launch.py} 提供这些入口：

\begin{longtable}{L{0.28\textwidth} L{0.18\textwidth} L{0.44\textwidth}}
\toprule
\textbf{Launch 参数} & \textbf{默认值} & \textbf{含义} \\
\midrule
\endhead
\param{params_file} & straight params & 传给 bridge/Nav2/控制器的 YAML 文件。\\
\param{network_interface} & \code{eth0} & Unitree SDK2 DDS 使用的网卡。外部电脑常见是 \code{enp5s0}。\\
\param{move_backend} & \code{nav2} & \code{nav2} 启动 Nav2 控制链；\code{cmd_vel} 只启动 bridge；采集脚本还支持 \code{direct}。\\
\param{use_bt_navigator} & \code{false} & 是否启动 BT navigator 和 behavior server。\\
\param{use_planner} & \code{false} & 是否启动 planner server。\\
\bottomrule
\end{longtable}

\section{minimal 参数文件}

\file{go2_nav2_minimal_controller_params.yaml} 适合无障碍直线采集。特点是：

\begin{itemize}
  \item \param{publish_point_cloud=false}，不依赖点云。
  \item \param{max_vy_mps=0.0}、\param{max_wz_radps=0.0}，限制侧向和旋转。
  \item local costmap 只有 inflation layer，不以点云避障为主。
  \item direct controller 使用 \param{distance_source=path_s}。
\end{itemize}

它更像“稳定采 RSSI/IMU 数据”的配置，而不是完整避障导航配置。

\section{straight/obstacle 参数文件}

\file{go2_nav2_straight_params.yaml} 适合 Nav2 obstacle 测试。特点是：

\begin{itemize}
  \item \param{publish_point_cloud=true}，发布 \topic{/go2/pointcloud}。
  \item local/global costmap 使用 \code{voxel_layer + inflation_layer}。
  \item controller 使用 DWB，并允许小幅 \code{vy} 和 \code{wz}。
  \item velocity smoother 限制速度变化。
  \item collision monitor 设置 FrontStop、SideStop、BodySlow、SideSlow 区域。
  \item \param{require_fresh_point_cloud_for_motion=true}，点云不新鲜时阻断运动。
\end{itemize}

\section{采集脚本如何编排 ROS 2}

\file{scripts/go2_collect_straight_line_nav2.sh} 是本项目的一键采集入口。它做的事情比普通 ROS 2 launch 多：

\begin{enumerate}
  \item 检查是否已有 Go2/Nav2 进程残留。
  \item source ROS 2 和本项目 overlay。
  \item 检查 \pkg{go2_nav2_bridge} 是否已经 build/source。
  \item 根据 \param{MOVE_BACKEND} 选择参数文件。
  \item 启动 RSSI 采集、IMU logger、Go2/Nav2 launch。
  \item 根据 backend 运行 direct controller、Nav2 client 或简单 \code{ros2 topic pub}。
  \item 停止后台进程，复制远程 RSSI 文件，写入 \file{run_metadata.txt}。
\end{enumerate}

推荐直线采集命令：

\begin{lstlisting}[language=bash]
cd /home/luping/桌面/RSSI/RSSI/code

RUN_DIR="$HOME/go2_rssi_runs/direct_1m_rssi" \
MOVE_BACKEND=direct \
RSSI_MODE=remote \
ROBOT_SSH=unitree@192.168.123.18 \
ROBOT_PROJECT_DIR=/home/unitree/rssi_go2 \
ROBOT_NETWORK_INTERFACE=eth0 \
BT_TTY=/dev/ttyACM1 \
DIRECT_DISTANCE_SOURCE=path_s \
MAX_RUNTIME_SEC=25 \
scripts/go2_collect_straight_line_nav2.sh enp5s0 1.0 0.20 1.0 0.5
\end{lstlisting}

输出目录通常包含：

\begin{itemize}
  \item \file{rssi_realtime_windows.csv}
  \item \file{rssi_raw_stream.csv}
  \item \file{imu_stream.csv} 或 \file{remote_imu_stream.csv}
  \item \file{nav2_launch.log}
  \item \file{run_metadata.txt}
\end{itemize}

\chapter{构建、运行和调试}
\label{chap:run-debug}

\section{从零构建}

外部电脑推荐 Ubuntu 22.04 + ROS 2 Humble：

\begin{lstlisting}[language=bash]
cd /home/luping/桌面/RSSI/RSSI/code/ros2_ws
source /opt/ros/humble/setup.bash
colcon build --packages-select go2_nav2_bridge
source install/setup.bash
\end{lstlisting}

检查：

\begin{lstlisting}[language=bash]
ros2 pkg prefix go2_nav2_bridge
ros2 pkg executables go2_nav2_bridge
\end{lstlisting}

应该能看到：

\begin{itemize}
  \item \code{go2_nav2_bridge}
  \item \code{go2_direct_straight_controller}
  \item \code{nav2_straight_path_client}
\end{itemize}

\section{只启动 bridge}

用于检查 Go2 DDS、odom、IMU、path\_s：

\begin{lstlisting}[language=bash]
ros2 run go2_nav2_bridge go2_nav2_bridge \
  --ros-args \
  --params-file /path/to/go2_nav2_minimal_controller_params.yaml \
  -p network_interface:=enp5s0
\end{lstlisting}

另一个终端检查：

\begin{lstlisting}[language=bash]
ros2 topic hz /odom
ros2 topic hz /go2/imu
ros2 topic echo /go2/path_s
ros2 run tf2_ros tf2_echo odom base_link
\end{lstlisting}

\section{手动发速度命令}

只在安全空旷环境做短时间测试：

\begin{lstlisting}[language=bash]
ros2 topic pub --rate 10 /cmd_vel geometry_msgs/msg/Twist \
  "{linear: {x: 0.05, y: 0.0, z: 0.0}, angular: {x: 0.0, y: 0.0, z: 0.0}}" \
  --times 20

ros2 topic pub --once /cmd_vel geometry_msgs/msg/Twist \
  "{linear: {x: 0.0, y: 0.0, z: 0.0}, angular: {x: 0.0, y: 0.0, z: 0.0}}"
\end{lstlisting}

\section{直线采集推荐路径}

推荐优先使用：

\begin{lstlisting}[language=bash]
MOVE_BACKEND=direct RSSI_MODE=off MAX_RUNTIME_SEC=20 \
scripts/go2_collect_straight_line_nav2.sh enp5s0 1.0 0.20 1.0 0.5
\end{lstlisting}

确认稳定后再打开 RSSI：

\begin{lstlisting}[language=bash]
MOVE_BACKEND=direct RSSI_MODE=remote \
ROBOT_SSH=unitree@192.168.123.18 \
scripts/go2_collect_straight_line_nav2.sh enp5s0 1.0 0.20 1.0 0.5
\end{lstlisting}

\section{Nav2 感知链路诊断}

在 obstacle/Nav2 测试前，先不运动，只检查点云和 costmap：

\begin{lstlisting}[language=bash]
RUN_DIR="$HOME/go2_rssi_runs/nav2_perception_diag" \
scripts/go2_diagnose_nav2_perception.sh enp5s0 12
\end{lstlisting}

诊断脚本会订阅：

\begin{itemize}
  \item \topic{/go2/pointcloud}
  \item \topic{/local_costmap/costmap}
\end{itemize}

并统计 front stop、side stop、body slow 等区域内点云数量。若点云 frame、外参或自过滤不对，先修这些，不要直接做运动测试。

\section{常用排查命令}

\begin{longtable}{L{0.34\textwidth} L{0.56\textwidth}}
\toprule
\textbf{命令} & \textbf{用途} \\
\midrule
\endhead
\code{ros2 node list} & 看节点是否启动。\\
\code{ros2 node info /go2_nav2_bridge} & 看 bridge 发布/订阅了哪些话题。\\
\code{ros2 topic list} & 看当前所有 topic。\\
\code{ros2 topic hz /go2/imu} & 看 IMU 是否持续发布。\\
\code{ros2 topic echo /go2/path_s} & 看累计路径进度是否变化。\\
\code{ros2 topic echo /cmd_vel} & 看最终进入 bridge 的速度命令。\\
\code{ros2 param get /go2_nav2_bridge max_vx_mps} & 查 bridge 速度限幅。\\
\code{ros2 lifecycle get /controller_server} & 查 Nav2 controller 是否 active。\\
\code{ros2 run tf2_ros tf2_echo odom base_link} & 查 TF 是否可用。\\
\code{tail -n 220 "$RUN_DIR/nav2_launch.log"} & 查采集 run 的 ROS 2/Nav2 日志。\\
\bottomrule
\end{longtable}

\chapter{ROS 2 和 RSSI/IMU 融合算法的关系}

\section{ROS 2 不直接等于融合算法}

本项目中，ROS 2 是数据和控制基础设施；RSSI/IMU 融合算法是定位研究主线。二者关系如下：

\begin{itemize}
  \item ROS 2 bridge 提供 \topic{/odom}、\topic{/go2/imu}、\topic{/go2/path_s}。
  \item RSSI 采集程序提供 \file{rssi_realtime_windows.csv}。
  \item 采集脚本把 RSSI、IMU、ROS 2 日志放到同一个 run directory。
  \item 融合算法按时间戳对齐 RSSI window 与 IMU/odom，输出预测路径进度 \(\hat{s}\)。
  \item Nav2 localization 输出不作为 RSSI 主算法输入，可以作为调试或对照参考。
\end{itemize}

\section{一条数据怎样进入训练或评估}

一次直线采集大致产生：

\begin{lstlisting}
run_dir/
  rssi_realtime_windows.csv
  rssi_raw_stream.csv
  imu_stream.csv
  remote_imu_stream.csv
  nav2_launch.log
  run_metadata.txt
\end{lstlisting}

融合算法使用它们时，直白流程是：

\begin{enumerate}
  \item RSSI window 给出当前无线指纹。
  \item IMU/odom 给出这个窗口附近的运动趋势和短时位移。
  \item 真实采集路线定义全局中心线 \(s\)。
  \item 训练时把 RSSI/IMU/odom 特征和真实 \(s\) 对齐。
  \item 在线时 RSSI 先猜位置，IMU/odom 限制跳变和短时连续性，最后输出 \(\hat{s}\)。
\end{enumerate}

\chapter{常见问题}

\section{ros2 run 找不到节点}

检查三件事：

\begin{lstlisting}[language=bash]
source /opt/ros/humble/setup.bash
source /home/luping/桌面/RSSI/RSSI/code/ros2_ws/install/setup.bash
ros2 pkg executables go2_nav2_bridge
\end{lstlisting}

如果没有 executables，重新 build。

\section{/odom 和 /go2/imu 没有数据}

可能原因：

\begin{itemize}
  \item \param{network_interface} 错了。
  \item Unitree DDS 网络没有连通。
  \item Unitree SDK2 动态库路径没有加入 \code{LD_LIBRARY_PATH}。
  \item bridge 没启动或启动后报错退出。
\end{itemize}

优先看：

\begin{lstlisting}[language=bash]
tail -n 120 "$RUN_DIR/nav2_launch.log"
ros2 node list
ros2 topic list
\end{lstlisting}

\section{Nav2 active 了但机器人不动}

按顺序查：

\begin{enumerate}
  \item \code{ros2 topic echo /cmd_vel_nav}：Nav2 controller 有没有输出。
  \item \code{ros2 topic echo /cmd_vel_smoothed}：velocity smoother 有没有输出。
  \item \code{ros2 topic echo /cmd_vel}：collision monitor 后还有没有最终速度。
  \item \code{ros2 param get /go2_nav2_bridge require_fresh_point_cloud_for_motion}：是否要求新鲜点云。
  \item \code{ros2 run tf2_ros tf2_echo odom base_link}：TF 是否正常。
\end{enumerate}

\section{点云导致机器人一直停}

如果 obstacle 参数中开启 \param{require_fresh_point_cloud_for_motion=true}，点云不更新或 frame 不对会阻断非零速度。先运行：

\begin{lstlisting}[language=bash]
scripts/go2_diagnose_nav2_perception.sh enp5s0 12
\end{lstlisting}

确认 \topic{/go2/pointcloud}、\topic{/local_costmap/costmap}、collision monitor 区域统计正常后，再做运动测试。

\section{什么时候用 direct，什么时候用 Nav2}

\begin{longtable}{L{0.24\textwidth} L{0.62\textwidth}}
\toprule
\textbf{模式} & \textbf{适合场景} \\
\midrule
\endhead
\code{MOVE_BACKEND=direct} & 推荐用于 RSSI/IMU 直线采集。链路短、变量少、结果更容易复现。\\
\code{MOVE_BACKEND=cmd_vel} & 最简单的手动/脚本速度测试，只适合短距离安全检查。\\
\code{MOVE_BACKEND=nav2} & 用于 obstacle、local planning、collision monitor 对照实验。调试成本更高，必须先确认点云和 TF 正常。\\
\bottomrule
\end{longtable}

\chapter{学习路线}

作为 ROS 2 新手，建议用本项目按以下顺序练习：

\begin{enumerate}
  \item 只 build \pkg{go2_nav2_bridge}，确认 \code{ros2 pkg executables} 能看到三个节点。
  \item 只启动 \rosnode{go2_nav2_bridge}，观察 \topic{/odom}、\topic{/go2/imu}、\topic{/go2/path_s}。
  \item 用 \code{ros2 topic pub} 发很短的 \topic{/cmd_vel}，确认 bridge 能控制 Go2。
  \item 运行 \rosnode{go2_direct_straight_controller}，理解 topic 闭环。
  \item 运行 \file{go2_collect_straight_line_nav2.sh}，理解 run directory 和采集文件。
  \item 运行 \file{go2_diagnose_nav2_perception.sh}，理解点云、costmap、QoS 和 TF。
  \item 最后再进入 Nav2 FollowPath、planner、BT navigator。
\end{enumerate}

\begin{notebox}{最重要的心智模型}
ROS 2 不是一个单独的大程序，而是一组节点通过标准接口组合起来。本项目里最稳的理解方式是沿着数据流看：Go2 DDS 进入 bridge，bridge 发布 ROS 2 标准话题，控制器或 Nav2 发布速度，bridge 把速度交给 Go2，采集脚本把 RSSI/IMU/odom/log 保存给融合算法使用。
\end{notebox}

\end{document}
